\begin{figure*}[t]
\centering
\begin{tikzpicture}[
  font=\small,
  >=stealth,
  node distance=2mm,
  graphbox/.style={draw, rounded corners=2pt, minimum width=24mm,
                   minimum height=22mm, align=center, inner sep=2pt,
                   fill=gray!5},
  blocklabel/.style={draw, rounded corners=1pt, fill=blue!10,
                     minimum width=18mm, minimum height=5.5mm,
                     font=\small\bfseries, align=center,
                     inner sep=1.5pt},
  termlabel/.style={font=\footnotesize\itshape, align=center,
                    text=blue!60!black, inner sep=1pt},
  arrow/.style={->, thick, >=stealth},
  fnode/.style={circle, draw, inner sep=0pt, minimum size=2.4mm,
                fill=white},
  enode/.style={circle, draw, inner sep=0pt, minimum size=2.4mm,
                fill=gray!30},
  lnode/.style={circle, draw, inner sep=0pt, minimum size=2.4mm,
                fill=black!70, text=white, font=\tiny},
  snode/.style={ellipse, draw, thick, inner sep=1pt,
                minimum width=10mm, minimum height=6mm,
                fill=orange!20},
  caplabel/.style={font=\footnotesize, align=center}
]

% =========================================================
% STAGE 1: Input attribution graph G
% =========================================================
\node[graphbox] (G) at (0,0) {};
\node[enode] (e1) at (-1.0, 0.7) {};
\node[enode] (e2) at (-1.0, 0.0) {};
\node[enode] (e3) at (-1.0,-0.7) {};
\node[fnode] (f1) at (-0.3, 0.85) {};
\node[fnode] (f2) at (-0.3, 0.3)  {};
\node[fnode] (f3) at (-0.3,-0.2)  {};
\node[fnode] (f4) at (-0.3,-0.75) {};
\node[fnode] (f5) at ( 0.3, 0.85) {};
\node[fnode] (f6) at ( 0.3, 0.3)  {};
\node[fnode] (f7) at ( 0.3,-0.2)  {};
\node[fnode] (f8) at ( 0.3,-0.75) {};
\node[lnode] (l1) at ( 1.0, 0.5) {};
\node[lnode] (l2) at ( 1.0,-0.3) {};
\foreach \s/\t in {e1/f1, e1/f2, e2/f2, e2/f3, e2/f6, e3/f4, e3/f7,
                   e3/f3, f1/f5, f2/f5, f2/f6, f3/f6, f3/f7, f4/f7,
                   f4/f8, f5/l1, f6/l1, f6/l2, f7/l2, f8/l2, f1/f6,
                   f2/f7, f3/f8}
  \draw[->, gray!50, line width=0.2pt] (\s) -- (\t);
\node[caplabel, below=1mm of G.south]
  {Input: $G=(V,E,W)$\\\scriptsize $|V|$$\sim$1000s, $|E|$$\sim$millions};

% =========================================================
% STAGE 2: Pruned graph G' after Block 1
% =========================================================
\node[graphbox, right=22mm of G] (Gp) {};
\node[enode] (pe1) at ($(Gp)+(-1.0, 0.5)$) {};
\node[enode] (pe2) at ($(Gp)+(-1.0,-0.4)$) {};
\node[fnode] (pf2) at ($(Gp)+(-0.3, 0.4)$) {};
\node[fnode] (pf3) at ($(Gp)+(-0.3,-0.2)$) {};
\node[fnode] (pf6) at ($(Gp)+( 0.3, 0.5)$) {};
\node[fnode] (pf7) at ($(Gp)+( 0.3,-0.3)$) {};
\node[lnode] (pl1) at ($(Gp)+( 1.0, 0.4)$) {};
\node[lnode] (pl2) at ($(Gp)+( 1.0,-0.4)$) {};
\foreach \s/\t in {pe1/pf2, pe2/pf3, pf2/pf6, pf3/pf6, pf3/pf7,
                   pf6/pl1, pf7/pl2, pf2/pf7}
  \draw[->, black, line width=0.4pt] (\s) -- (\t);
\node[caplabel, below=1mm of Gp.south]
  {Pruned $(V',E')$\\\scriptsize keeps top $C^V,C^E$};

% =========================================================
% STAGE 3: Clustered graph after Block 2 (cycle present)
% =========================================================
\node[graphbox, right=22mm of Gp] (Gc) {};

% Supernodes
\node[snode, minimum width=7mm, minimum height=5mm]
  (s1) at ($(Gc)+(-0.55, 0.5)$) {};
\node[snode, minimum width=7mm, minimum height=5mm]
  (s2) at ($(Gc)+(-0.55,-0.5)$) {};
\node[snode, minimum width=7mm, minimum height=5mm]
  (s3) at ($(Gc)+( 0.55, 0.5)$) {};
\node[snode, minimum width=7mm, minimum height=5mm]
  (s4) at ($(Gc)+( 0.55,-0.5)$) {};

% Embedding/logit columns (named so we can draw edges to/from them)
\node[enode] (ec1) at ($(Gc)+(-1.15, 0.5)$) {};
\node[enode] (ec2) at ($(Gc)+(-1.15,-0.5)$) {};
\node[lnode] (lc1) at ($(Gc)+( 1.15, 0.4)$) {};
\node[lnode] (lc2) at ($(Gc)+( 1.15,-0.4)$) {};

% Feature dots inside each supernode
\foreach \dx/\dy in {-0.65/0.5, -0.45/0.5}
  \fill[black] ($(Gc)+(\dx,\dy)$) circle (0.5mm);
\foreach \dx/\dy in {-0.65/-0.5, -0.45/-0.5}
  \fill[black] ($(Gc)+(\dx,\dy)$) circle (0.5mm);
\foreach \dx/\dy in {0.45/0.5, 0.65/0.5}
  \fill[black] ($(Gc)+(\dx,\dy)$) circle (0.5mm);
\foreach \dx/\dy in {0.45/-0.5, 0.65/-0.5}
  \fill[black] ($(Gc)+(\dx,\dy)$) circle (0.5mm);

% Boundary edges: embeddings -> left supernodes, right supernodes -> logits
\draw[->, thick] (ec1.east) -- (s1.west);
\draw[->, thick] (ec2.east) -- (s2.west);
\draw[->, thick] (s3.east)  -- (lc1.west);
\draw[->, thick] (s4.east)  -- (lc2.west);

% Inter-supernode directed edges
\draw[->, thick] (s1.east)  -- (s3.west);
\draw[->, thick] (s2.east)  -- (s4.west);
\draw[->, thick] (s1.south) -- (s2.north);
\draw[->, thick] (s3.south) -- (s4.north);

% Highlighted 2-cycle (red): one direction solid, opposite dashed
\draw[->, thick, red!80!black]
  (s3.south west) to[bend left=20] (s2.north east);
\draw[->, thick, red!80!black, dashed]
  (s2.north east) to[bend left=20] (s3.south west);

\node[caplabel, below=1mm of Gc.south]
  {Clustered $\varphi$\\\scriptsize Ncut on $\tilde S$ (cycle!)};

% =========================================================
% STAGE 4: Final supernode graph (DAG)
% =========================================================
\node[graphbox, right=22mm of Gc] (Gf) {};

\node[snode, minimum width=7mm, minimum height=5mm]
  (fs1) at ($(Gf)+(-0.55, 0.5)$) {};
\node[snode, minimum width=7mm, minimum height=5mm]
  (fs2) at ($(Gf)+(-0.55,-0.5)$) {};
\node[snode, minimum width=7mm, minimum height=5mm]
  (fs3) at ($(Gf)+( 0.55, 0.5)$) {};
\node[snode, minimum width=7mm, minimum height=5mm]
  (fs4) at ($(Gf)+( 0.55,-0.5)$) {};

\node[enode] (ef1) at ($(Gf)+(-1.15, 0.5)$) {};
\node[enode] (ef2) at ($(Gf)+(-1.15,-0.5)$) {};
\node[lnode] (lf1) at ($(Gf)+( 1.15, 0.4)$) {};
\node[lnode] (lf2) at ($(Gf)+( 1.15,-0.4)$) {};

\foreach \dx/\dy in {-0.65/0.5, -0.45/0.5}
  \fill[black] ($(Gf)+(\dx,\dy)$) circle (0.5mm);
\foreach \dx/\dy in {-0.65/-0.5, -0.45/-0.5}
  \fill[black] ($(Gf)+(\dx,\dy)$) circle (0.5mm);
\foreach \dx/\dy in {0.45/0.5, 0.65/0.5}
  \fill[black] ($(Gf)+(\dx,\dy)$) circle (0.5mm);
\foreach \dx/\dy in {0.45/-0.5, 0.65/-0.5}
  \fill[black] ($(Gf)+(\dx,\dy)$) circle (0.5mm);

% Boundary edges
\draw[->, thick] (ef1.east) -- (fs1.west);
\draw[->, thick] (ef2.east) -- (fs2.west);
\draw[->, thick] (fs3.east) -- (lf1.west);
\draw[->, thick] (fs4.east) -- (lf2.west);

% Inter-supernode edges
\draw[->, thick] (fs1.east)  -- (fs3.west);
\draw[->, thick] (fs2.east)  -- (fs4.west);
\draw[->, thick] (fs1.south) -- (fs2.north);
\draw[->, thick] (fs3.south) -- (fs4.north);

% Resolved cycle: only one direction retained
\draw[->, thick, darkgreen]
  (fs3.south west) to[bend left=20] (fs2.north east);

\node[caplabel, below=1mm of Gf.south]
  {Supernode graph $G_{SN}$\\\scriptsize DAG};
% =========================================================
% PIPELINE ARROWS (drawn first; sit below the labels visually)
% =========================================================
\draw[arrow] (G.east)  -- (Gp.west);
\draw[arrow] (Gp.east) -- (Gc.west);
\draw[arrow] (Gc.east) -- (Gf.west);

% =========================================================
% BLOCK + TERM LABELS, stacked ABOVE the graph boxes so they
% don't overlap arrows or boxes.
% Graph boxes top ~ y = +1.1; place stack starting at y = +1.7.
% =========================================================
\node[blocklabel] (b1) at ($(G.east)!0.5!(Gp.west) + (0, 1.7)$)
  {Block 1\\\scriptsize Prune};
\node[termlabel, above=0.8mm of b1]
  {term (A): $C^V, C^E$ under (F3)};

\node[blocklabel] (b2) at ($(Gp.east)!0.5!(Gc.west) + (0, 1.7)$)
  {Block 2\\\scriptsize Spectral};
\node[termlabel, above=0.8mm of b2]
  {term (B): Ncut($\tilde S$) $+$ (F2)};

\node[blocklabel] (b3) at ($(Gc.east)!0.5!(Gf.west) + (0, 1.7)$)
  {Block 3\\\scriptsize DAG proj.};
\node[termlabel, above=0.8mm of b3]
  {constraint (F1)};

% Tiny vertical tick from each block label down to its arrow,
% so the reader sees the label "applies to" that arrow.
\draw[dotted, gray] (b1.south) -- ($(G.east)!0.5!(Gp.west)$);
\draw[dotted, gray] (b2.south) -- ($(Gp.east)!0.5!(Gc.west)$);
\draw[dotted, gray] (b3.south) -- ($(Gc.east)!0.5!(Gf.west)$);

% =========================================================
% LEGEND at bottom
% =========================================================
\node[anchor=west, font=\scriptsize] at ($(G.south west)+(0,-1.5)$)
  {\textbf{Legend:}};
\node[anchor=west, font=\scriptsize] at ($(G.south west)+(1.2,-1.5)$)
  {\tikz\node[enode, scale=0.9]{};\, embedding};
\node[anchor=west, font=\scriptsize] at ($(G.south west)+(3.5,-1.5)$)
  {\tikz\node[fnode, scale=0.9]{};\, feature};
\node[anchor=west, font=\scriptsize] at ($(G.south west)+(5.4,-1.5)$)
  {\tikz\node[lnode, scale=0.9]{};\, logit};
\node[anchor=west, font=\scriptsize] at ($(G.south west)+(7.2,-1.5)$)
  {\tikz\node[snode, scale=0.6]{};\, supernode};
\node[anchor=west, font=\scriptsize, text=red!80!black]
  at ($(G.south west)+(10.2,-1.5)$) {cycle to resolve};
\node[anchor=west, font=\scriptsize, text=darkgreen]
  at ($(G.south west)+(13.0,-1.5)$) {kept direction};

\end{tikzpicture}
\caption{The ASG pipeline as block-coordinate descent on
$\mathcal{J}_K$. \textbf{Block~1} prunes $(V', E')$ by
cumulative-mass thresholding on the combined evidence scores
$C^V, C^E$, solving term~(A) under the budget constraint (F3) as a
unit-cost knapsack. \textbf{Block~2} partitions the surviving
features by spectral clustering on the bidirectional affinity
$\tilde S$, solving term~(B) at fixed $(V', E')$ via the
Shi--Malik relaxation and absorbing layer locality (F2) through the
exponential decay in $\tilde S$. \textbf{Block~3} restores the DAG
constraint (F1) by per-pair direction selection (eliminating the
2-cycle highlighted in red) followed by SCC-split-at-widest-span.
The coupling between blocks is one-directional within a sweep
($(V', E') \to \tilde S \to \varphi$), so the single forward sweep
is a fixed point of the coordinate iteration. Granularity $K$ is
selected externally over the grid $\mathcal{K}$ (not shown).}
\label{fig:pipeline}
\end{figure*}